\section{Two other reporting rules}\label{sec:alternatives}

The report in Section~\ref{sec:observation}, which we call the
\emph{refined report}, keeps every leaf value when all the leaves pass
their tests. We now ask what happens if that information is replaced by one
mark. The resulting loss can be computed exactly and is smaller than
the Gaussian gain. A second rule uses one additional independent input:
it reveals the original inputs off the exceptional event and the extra
input on that event. The first modification shows that the all-pass
leaf values are dispensable; the second isolates the conditional-variance
saving on the exceptional event.

\subsection{Merging the all-leaves-pass reports}\label{subsec:coarse-mark}

Start with any integer $m\ge3$, intervals $I_1,\ldots,I_m$, and a
selector $\iota:\R\to\{1,\ldots,m\}$ with finitely many cut points.
Fix all endpoint and tie conventions. As before, on inputs
$(u,z_1,\ldots,z_m)$ put
\[
 L=u+\sum_{j=1}^m z_j,\qquad
 B=\{z_j\in I_j\text{ for every }j\},
\]
\[
 P=\{z_j\in I_j\text{ for }j\ne\iota(u)\},\qquad
 A=P\setminus B.
\]
Define the \emph{coarse report} $\mathcal O$ by
\[
\begin{array}{c|l}
 \text{Event}&\text{Information reported}\\ \hline
 A&\text{the mark }X\\
 B&\text{the mark }Y\\
 P^c&(O,u,\iota(u),(z_j)_{j\ne\iota(u)}).
\end{array}
\]
The marks $X,Y$ and the ordinary-report tag $O$ are distinct. The rule
differs from the refined report only on $B$. The next lemma separates
its pointwise cost from the choice of an input law, so that other
interval constructions can use the same reporting rule.

\begin{lemma}[Cost and error of the coarse report]\label{lem:coarse-report}
On every product of finite input alphabets, $\mathcal O$ has finite
range, and each output indicator is a finite linear combination of
functions of at most $m$ input coordinates.

Let the inputs be independent, centered, and of variance one; their
marginal laws need not agree. For any fixed constants $b_0,b_1\in\R$,
define
\begin{equation}\label{eq:coarse-deficit}
 g_{b_0,b_1}
 =\one_A\bigl(1-(L-b_0)^2\bigr)
   +\one_B\bigl(1-(L-b_1)^2\bigr).
\end{equation}
Then the optimal residual error
$\mathcal R=\E[(L-\E[L\mid\mathcal O])^2]$ satisfies
\begin{equation}\label{eq:coarse-residual}
 1-\mathcal R\ge\E g_{b_0,b_1},\qquad
 \Var(\E[L\mid\mathcal O])\ge m+\E g_{b_0,b_1}.
\end{equation}
\end{lemma}

\begin{proof}
The $X$ indicator has the pointwise expansion \eqref{eq:indicator}.
The $Y$ indicator depends only on the $m$ leaves. An ordinary output
specifies $u$, its selected index $i$, and the $m-1$ unselected leaf
values. If these data are attained, one of the reported leaves fails
its interval test. They therefore force $P^c$ for every value of the
omitted leaf $z_i$, so the indicator uses just the $m$ reported inputs.
An inconsistent or unattained output label has zero indicator. Thus
the assertion holds pointwise, including endpoints and selector ties.

Predict $L$ by $b_0$ on $X$, by $b_1$ on $Y$, and by
$u+\sum_{j\ne\iota(u)}z_j$ on an ordinary report. For each fixed
$i$, the event $E_i=P^c\cap\{\iota(u)=i\}$ and all its reported data
depend only on $u$ and the leaves other than $z_i$. Independence gives
\[
 \E[\one_{E_i}z_i]=0,\qquad
 \E[\one_{E_i}z_i^2]=\Pp(E_i).
\]
The ordinary reports consequently contribute exactly $\Pp(P^c)$ to
the squared prediction error. The total error of this predictor is
\[
 \Pp(P^c)+\E[\one_A(L-b_0)^2]+\E[\one_B(L-b_1)^2]
 =1-\E g_{b_0,b_1}.
\]
Conditional expectation minimizes squared error, proving the first
inequality in \eqref{eq:coarse-residual}. The second follows because
$L$ is centered with variance $m+1$.
\end{proof}

Now fix the Gaussian intervals and selector from
Lemma~\ref{lem:gaussian}. Write
\[
 M_\mu=\sum_{j=1}^m\mu_j,\qquad b=M_\mu-1,\qquad
 p_*=\Pp(B)>0,\qquad S=\sum_{j=1}^m s_j^2.
\]
The Gaussian calculation gave
$d_A=\E[\one_A(1-(L-b)^2)]=p_*(1+S-J)$ with $J<1/4$.
In Lemma~\ref{lem:coarse-report} take $b_0=b$ and $b_1=M_\mu$.
On $B$, the central Gaussian remains unrestricted and the leaves are
independently restricted to their intervals. Thus $L-M_\mu$ has
conditional mean zero and variance $1+S$, giving
\begin{equation}\label{eq:coarse-gaussian-gain}
 \E g_{b,M_\mu}=d_A-p_*S=p_*(1-J)>\frac34p_*.
\end{equation}
The term $p_*S$ is precisely the added error on $B$: the refined
predictor had error $u$ there, while merging its reports leaves the
conditional leaf variance $S$ unobserved. Equation
\eqref{eq:coarse-gaussian-gain} is an identity for this predictor's
gain; \eqref{eq:coarse-residual} is the resulting lower bound for
optimal retained variance.

All these Gaussian choices are now fixed. Let centered variance-one
laws $Z_n$ converge weakly to a standard normal with uniformly bounded
fourth moments, and take $m+1$ independent copies as inputs. The fixed
function $g_{b,M_\mu}$ has quadratic growth; its discontinuities lie
on the finitely many interval and selector hyperplanes.
Lemma~\ref{lem:transfer} and \eqref{eq:coarse-gaussian-gain} imply
that eventually $\E g_{b,M_\mu}>\eta$, where $\eta=p_*/2>0$.
Lemma~\ref{lem:coarse-report} then gives retained variance greater
than $m+\eta$ and the required pointwise cell cost. Therefore
Proposition~\ref{prop:amplification} applies with
\[
 (r,d,\eta)=(m+1,m,p_*/2).
\]
This criterion and Lemma~\ref{lem:readout} give
Theorem~\ref{thm:main} by the coarse rule as well.

\subsection{Switching to an independent score}\label{subsec:extra-block}

Keep the fixed Gaussian event $A$ and predictor $b$ above, and put
$q=m+1$. Add one independent centered variance-one input $w$ to the
original tuple $\mathbf z=(u,z_1,\ldots,z_m)$. Define
\[
 \mathcal O_+(\mathbf z,w)=
 \begin{cases}
   (\mathrm{off},\mathbf z),&\mathbf z\notin A,\\
   (\mathrm{on},w),&\mathbf z\in A.
 \end{cases}
\]
Off $A$ the original sum $L$ is known and $w$ is omitted. On $A$
the extra input is known and the original sum must be predicted.
The rule gains variance whenever the conditional variance of $L$ on
$A$ is smaller than the one unit of residual variance of $w$.

This rule has $q+1$ inputs but cell cost at most $q$. Indeed, an
off-$A$ output fixes the first $q$ inputs and places no condition on
$w$. The indicator of an on-$A$ output $(\mathrm{on},w_0)$ is
\[
 \one_A(\mathbf z)\one_{\{w=w_0\}}.
\]
Multiplying \eqref{eq:indicator} by the last value test expresses
this indicator as a sum of functions of at most $m+1=q$ inputs.
The tags keep the two types of report separate, and unattained labels
have zero indicators. These pointwise expansions also cover the
prescribed endpoints and ties. The rule has finite range on every
finite product alphabet.

To establish a gain for the approximating laws, first transfer the
unconditional deficit, before conditioning on $A$. For independent
copies of any centered variance-one law $Z_n$ tending to a standard
normal with uniformly bounded fourth moments, Lemma~\ref{lem:transfer} gives
\[
 \Delta_{b,n}:=\E[\one_A(1-(L-b)^2)]
 \longrightarrow d_A>\frac34p_*.
\]
Thus, for all sufficiently large $n$, $\Delta_{b,n}>\eta=p_*/2$.
Since $\Delta_{b,n}\le\Pp(A)$, these laws have positive probability
of $A$. Only now fix such an $n$, write $\Delta_b=\Delta_{b,n}$, and define
\[
 p=\Pp(A)>0,\qquad c=\E[L\mid A].
\]
Independence of $w$ from the original tuple shows that
$\E[L+w\mid\mathcal O_+]$ equals $L$ off $A$ and $c+w$ on $A$.
The sum $L+w$ is centered, so its conditional mean is centered too.
Using $\E w=0$, $\E w^2=1$, and $\E L^2=q$ gives the exact identity
\begin{align}
 \Var(\E[L+w\mid\mathcal O_+])
 &=\E[\one_{A^c}L^2]+p(c^2+1)\notag\\
 &=q+p\bigl(1-\Var(L\mid A)\bigr).
 \label{eq:switching-variance}
\end{align}
The same input law satisfies
\begin{equation}\label{eq:switching-correction}
 p\bigl(1-\Var(L\mid A)\bigr)
 =\Delta_b+p(c-b)^2>\eta.
\end{equation}
The identities in \eqref{eq:switching-variance} and the equality in
\eqref{eq:switching-correction} hold more generally for any independent
centered variance-one original inputs and an independent centered
variance-one $w$, whenever $p>0$, with
$\Delta_b=\E[\one_A(1-(L-b)^2)]$. Their marginal laws need not agree.
The nonnegative correction $p(c-b)^2$ records the improvement from
the fixed Gaussian predictor $b$ to the actual conditional mean $c$.
For the fixed approximating law, \eqref{eq:switching-variance}
therefore exceeds $q+\eta$.
This argument uses convergence of one fixed predictor error, not
continuity of conditional variances.

The pointwise cost and the transferred gain verify
Proposition~\ref{prop:amplification} with the distinct triple
\[
 (r,d,\eta)=(m+2,m+1,p_*/2).
\]
The criterion again proves Theorem~\ref{thm:main}, with ratio
multiplier at least $1+\eta/(m+1)>1$. For both rules in this section,
the Gaussian parameters and $\eta$ stay fixed throughout. Each stage
starts from one fixed finite observation and chooses its own finite
averaging count; the fourth-moment bound is uniform in that count,
not across all stages.
